\documentclass[11pt]{article}
\usepackage[margin=1in]{geometry}
\usepackage{setspace}
\usepackage{lineno}
\usepackage{amsmath,amssymb}
\usepackage{booktabs}
\usepackage{threeparttable}
\usepackage{threeparttablex}
\usepackage{longtable}
\usepackage{tabularx}
\usepackage{array}
\usepackage{graphicx}
\usepackage{caption}
\usepackage{subcaption}
\usepackage{float}
\usepackage{placeins}
\usepackage{microtype}
\usepackage{xcolor}
\usepackage{url}
\usepackage{csquotes}
\usepackage[
  backend=biber,
  style=authoryear,
  natbib=true,
  maxcitenames=2,
  maxbibnames=20,
  giveninits=true,
  doi=true,
  url=true,
  isbn=false,
  hyperref=true,
  backref=true
]{biblatex}
\usepackage[
  hidelinks
]{hyperref}
\usepackage{enumitem}
\usepackage{titlesec}
\usepackage{fancyhdr}
\usepackage{lastpage}
\addbibresource{references_final.bib}
\hypersetup{
  pdftitle={Do Remittances Buffer Household Shocks? Evidence on Food Insecurity in Kyrgyzstan and Uzbekistan},
  pdfauthor={Foziljon Alisherov and Mukhayyo Djuraeva},
  pdfsubject={Remittances, household shocks, and food insecurity in Central Asia},
  pdfkeywords={remittances; household shocks; food insecurity; resilience; migration; Kyrgyzstan; Uzbekistan; Central Asia}
}
\pagestyle{fancy}
\fancyhf{}
\rhead{\thepage}
\lhead{Remittances, shocks, and food insecurity}
\captionsetup{font=small,labelfont=bf}
\renewcommand{\arraystretch}{1.15}
\doublespacing
\linenumbers


\begin{document}

\begin{titlepage}
\centering
{\Large\bfseries\textcolor{teal!45!black}{Do Remittances Buffer Household Shocks?\\Evidence on Food Insecurity in Kyrgyzstan and Uzbekistan}\par}
\vspace{0.4em}
\textcolor{teal!45!black}{\rule{0.94\textwidth}{0.6pt}}

\vspace{1.0em}
{\large\textit{Evidence from the Life in Kyrgyzstan Study and Listening to the Citizens of Uzbekistan Survey}\par}

\vspace{1.3em}
{\large\bfseries Working Paper\par}

\vspace{1.3em}
{\large\bfseries Foziljon Alisherov and Mukhayyo Djuraeva\par}

\vspace{1.1em}
{\large Research Assistant, New Uzbekistan University; New Uzbekistan University\par}

\vspace{0.8em}
{\large ORCID: \href{https://orcid.org/0009-0004-9451-0518}{0009-0004-9451-0518}; \href{https://orcid.org/0000-0001-6163-7513}{0000-0001-6163-7513}\par}

\vspace{0.8em}
{\large Emails: \href{mailto:f.alisherov@newuu.uz}{f.alisherov@newuu.uz}; \href{mailto:m.djuraeva@newuu.uz}{m.djuraeva@newuu.uz}\par}

\vspace{0.8em}
{\large University profile: \href{https://www.newuu.uz/en/employee/view/faculty-members/419}{New Uzbekistan University profile}\par}

\vspace{1.1em}
{\large August 2026\par}

\vspace{1.1em}
{\large Corresponding author: Foziljon Alisherov\par}
{\large Email: \href{mailto:f.alisherov@newuu.uz}{f.alisherov@newuu.uz}\par}
\end{titlepage}

\begin{abstract}
This paper asks whether the positive association between household shocks and food insecurity is weaker among remittance-receiving households in Kyrgyzstan and Uzbekistan. The analysis uses the 2019 wave of the Life in Kyrgyzstan Study and Listening to the Citizens of Uzbekistan rounds 49-82. Because the surveys differ in design, unit of observation, reference period and shock measurement, all models are country-specific and observational rather than pooled. Kazakhstan is included only as regional benchmark context. The preferred Kyrgyzstan model estimates a negative remittance-shock interaction of -0.2140 (95\% CI -0.6549 to 0.2269; p=0.3415), which is directionally consistent with a weaker shock-food-insecurity association but imprecise. The preferred Uzbekistan broad-shock model estimates a negative interaction of -0.5406 (95\% CI -1.0415 to -0.0398; p=0.03437). Four-group predictions are used to interpret the interaction on the food-insecurity raw-score scale and to separate baseline remittance differences from shock-period differences. The Uzbekistan household fixed-effects estimate remains negative but is smaller and imprecise (-0.1771; 95\% CI -0.5515 to 0.1973; p=0.3539). L2CU estimates are unweighted because the available \texttt{popw} documentation was not approved for this analysis. The findings are compatible with remittance-related moderation of household shock vulnerability, especially in Uzbekistan, but the evidence remains associational and subject to selection, timing and measurement limitations.
\end{abstract}

\noindent\textbf{Keywords:} remittances; household shocks; food insecurity; resilience; migration; Kyrgyzstan; Uzbekistan; Central Asia

\clearpage

\section{Introduction}

Household shocks are a central source of food-security risk. A job loss, illness, death, agricultural loss or climate-related event can reduce income, raise expenses and disrupt the household's ability to obtain food. These events are often not isolated from other household decisions. Families may borrow, sell assets, change labor supply, reduce non-food spending or draw on help from relatives. In migration-dependent settings, remittances are one possible source of support. They may add liquidity when a household faces stress, but they also identify households with migration networks, different labor-market histories and potentially different underlying resources.

This paper asks a focused research question: is the positive association between household shocks and food insecurity weaker among remittance-receiving households in Kyrgyzstan and Uzbekistan? The outcome is a higher-worse food-insecurity raw score. The central empirical test is therefore whether the shock-associated increase in that score is smaller for households that receive remittances. A negative remittance-shock interaction is consistent with such a weaker adverse association, but it is not interpreted as a causal estimate.

The question is important for Central Asia because labor migration and remittances are embedded in household livelihood strategies. The World Bank's migration-development synthesis emphasizes that origin-country households may receive remittances and knowledge transfers while also facing the absence and risks associated with migration \parencite{WorldBank2023}. Classic work on migration and risk argues that households may use migration to diversify income sources \parencite{StarkLevhari1982}, while remittance-motivation research shows that transfers may reflect family arrangements, risk sharing and mutual interest rather than a single simple motive \parencite{LucasStark1985}. These ideas are relevant to resilience research, but they do not by themselves establish what happens in a particular household survey when food-insecurity outcomes and recent shocks are observed together.

The empirical design uses two main datasets. For Kyrgyzstan, the analysis uses the 2019 wave of the Life in Kyrgyzstan Study, a multi-topic longitudinal survey with household and individual modules \parencite{LIK2023}. For Uzbekistan, it uses Listening to the Citizens of Uzbekistan, a high-frequency World Bank phone survey that tracks household welfare, food security, employment, remittances, services and shocks \parencite{WorldBank2025a,WorldBank2025b}. Kazakhstan is not a third country in the mechanism analysis. Its FIES data are used only for regional food-insecurity and demographic benchmark context because verified remittance and household-shock mechanism variables are not available.

The paper makes three contributions. First, it applies the same conceptual moderation question to two Central Asian settings without pooling incompatible respondent records. Second, it reports four-group adjusted predictions and interaction contrasts rather than relying only on the coefficient of the product term. Third, it distinguishes primary evidence, robustness qualifications and benchmark context. Kyrgyzstan is classified as directional but imprecise; Uzbekistan is classified as a moderate conditional association with limitations; Uzbekistan household fixed effects are reported as a prominent qualification; and the work-loss-only Uzbekistan result is kept as secondary event-specific exploratory evidence.

The preview of findings is restrained. In Kyrgyzstan, the preferred interaction is negative but imprecise. In Uzbekistan, the preferred broad-shock interaction is negative and more precise, but the remittance-plus-shock cell is modest and the household fixed-effects estimate is attenuated and imprecise. The results are best read as evidence that remittance status conditions the observed shock-food-insecurity association in these data, especially in Uzbekistan, while leaving open selection, timing and measurement explanations.

\section{Literature review and research gap}

\subsection{Household shocks and food insecurity}

Food insecurity reflects constrained access to adequate food and is closely related to household resources, prices, labor-market conditions and coping capacity. The Food Insecurity Experience Scale is an experience-based measure built from eight questions about food-related hardship \parencite{FAO2026}. \textcite{Cafiero2018} describe the FIES measurement model and the role of calibration for comparable population estimates. This paper uses raw FIES-style scores rather than official calibrated estimates, so the literature supports the interpretation of the outcome as reported food-access hardship while also cautioning against over-comparison across surveys.

Empirical work links shocks to food-security and welfare outcomes through multiple channels. \textcite{Hoddinott2006} emphasizes that household shocks can shape consumption, assets and poverty dynamics within and across households. \textcite{AkterBasher2014} show that food-price and income shocks were associated with worse food security in rural Bangladesh, especially for poorer households. \textcite{Ansah2021} extend the discussion to interacting climate, health, pest and price shocks and show that coping strategies matter for food-security outcomes. These studies support the expectation that household shocks are positively associated with higher food-insecurity scores.

\subsection{Migration and remittances as informal insurance}

The remittance literature offers a plausible reason to expect moderation. \textcite{StarkLevhari1982} frame migration as a household response to risk. \textcite{LucasStark1985} show that remittance motives can include mutual family arrangements and risk sharing. \textcite{ChoiYang2007} provide influential evidence from the Philippines that remittances may respond to negative income shocks and that consumption is smoother among households with migrants. \textcite{CombesEbeke2011} similarly connect remittances to lower household consumption instability in developing countries. Together, this literature supports the idea that remittances can be associated with consumption smoothing, while also emphasizing that the empirical setting matters.

The policy literature reaches a similar boundary from a different direction. Shock-responsive and adaptive social protection frameworks focus on public systems that can respond to household and covariate shocks \parencite{OBrien2018,Bowen2020}. Private remittances can be part of a household's resource envelope, but they are not public entitlements and are not guaranteed. The present study therefore connects remittance heterogeneity to vulnerability analysis without treating private transfers as substitutes for formal social protection.

\subsection{Selection, endogeneity and alternative explanations}

Remittance receipt is not random. Migrant-sending households may have more resources, better networks or different risk exposure before a shock occurs. Transfers may also be reactive: relatives may send money after illness, job loss or another event. \textcite{AzamGubert2006} review evidence showing that remittance motives and household selection are complex. \textcite{Kakhkharov2021} address selection in their Uzbekistan remittance-expenditure study, and \textcite{Wang2021} explicitly treat endogeneity and heterogeneity in their Kyrgyzstan remittance-spending analysis. These concerns are central to the interpretation here.

Because of these concerns, the paper does not ask whether remittances causally reduce food insecurity. It asks whether the observed shock-food-insecurity association differs by remittance status. This distinction is more than wording. It determines model interpretation, control selection and policy language. Controls should not mechanically include variables such as current income, asset sales or coping behavior when those variables may lie downstream of shocks or remittances. Household fixed effects can help assess stable unobserved differences in the Uzbekistan panel, but they cannot resolve timing, measurement error or all time-varying confounding.

\subsection{Evidence from Kyrgyzstan}

Kyrgyzstan is especially relevant because LiK provides unusually rich household and individual microdata for Central Asia. \textcite{Brueck2014} describe the limited availability of longitudinal microdata in the region and present LiK as a major panel resource. The official LiK data record describes waves from 2010 through 2019, coverage across Kyrgyz regions and modules on migration, expenditure, shocks, agriculture and individual welfare \parencite{LIK2023}. This survey structure makes Kyrgyzstan a strong setting for a household resilience question.

Existing LiK-based research also informs the present design. \textcite{Wang2021} use LiK to examine remittances and household spending strategies and find that remittances are related to household budget allocation, with careful attention to fixed effects and heterogeneity. That work does not answer the present food-insecurity moderation question, but it supports the relevance of remittances, household spending and selection concerns in Kyrgyzstan.

\subsection{Evidence from Uzbekistan}

Uzbekistan has a thinner peer-reviewed evidence base on the precise remittance-shock-food-insecurity nexus. \textcite{Kakhkharov2021} study South-South migration, remittances and household expenditures in Uzbekistan and use instrumental-variable regressions to address selection. Their findings show that remittance-receiving households differ in expenditure allocation, which supports the relevance of remittances for household welfare analysis in Uzbekistan. The World Bank's L2CU documentation adds a strong survey foundation: L2CU monitors household income, coping, food security, employment, remittances, services and shocks through repeated phone interviews \parencite{WorldBank2025a,WorldBank2025b,Uochi2025}.

The remaining gap is specific. Verified literature supports remittance endogeneity, consumption smoothing, FIES measurement, household shocks and the L2CU data source. It does not fully settle whether remittances condition the relationship between verified household shocks and FIES-style food insecurity in Uzbekistan. This paper addresses that gap with country-specific observational models and by reporting the fixed-effects qualification prominently.

The literature also clarifies why the dependent variable should be interpreted carefully. FIES-style items capture experiences of constrained access to food, not the full multidimensional food-security concept. Availability, access, utilization, stability, agency and sustainability can all matter, but the eight experience items focus most directly on access. This makes the measure well suited to household-level welfare analysis, especially when shocks affect purchasing power or coping behavior. It also means the paper should not claim to measure nutrition, dietary quality or agricultural production directly. The outcome records food-related hardship experiences.

The household-shock literature further suggests that shocks differ in both immediacy and pathway. Employment shocks may reduce cash income quickly. Health shocks may combine lost labor time with new expenses. Agricultural and climate events may affect own production, livestock, input costs or local prices. Death in the household can have emotional, economic and caregiving dimensions. A broad shock indicator is therefore a simplification, but it is useful for testing whether the overall relationship between adverse events and food insecurity varies by remittance status.

The remittance literature provides a useful but incomplete bridge between household risk and food insecurity. Informal-insurance models imply that migrants and origin households may share risk across locations. If local shocks are not perfectly correlated with migrant earnings, transfers can smooth resources. Yet the same literature warns that remittance flows depend on migrant employment, destination-country conditions, household bargaining and expectations about future support. A remittance indicator therefore captures both resources and relationships. This is why the paper interprets remittance receipt as a conditioning variable rather than a treatment.

Central Asian evidence reinforces that caution. Kyrgyzstan is often discussed as highly remittance dependent, but household-level studies show that remittances do not translate mechanically into a single welfare pattern. Uzbekistan research similarly shows that remittance-receiving households differ in expenditure allocation and that selection is analytically important. These findings justify the paper's focus on heterogeneity, but they also limit the language used in the discussion. The manuscript can say that the observed shock gradient differs by remittance status; it cannot say that remittances are sufficient for resilience.

\section{Conceptual framework and hypotheses}

The conceptual framework has three components: household shocks, food insecurity and remittance receipt. Shocks can worsen food access by reducing cash income, increasing medical or funeral expenses, disrupting agricultural production, or reducing labor availability. Food insecurity is observed as a higher raw score on eight food-access hardship items. Remittance receipt may be associated with additional liquidity or support from migrant networks, but it may also be correlated with unobserved household characteristics.

The first hypothesis is that household shocks are positively associated with food insecurity among households without remittance receipt. The second is that the shock-food-insecurity association is weaker among remittance-receiving households. The empirical marker for the second hypothesis is a negative coefficient on the interaction between remittance receipt and shock exposure, because the outcome is higher-worse. The third expectation is that results will differ in precision and robustness across countries because the surveys have different designs, recall periods and shock definitions.

The mechanism is intentionally described as plausible rather than demonstrated. Remittances may help smooth consumption, but transfers may also be sent because hardship has already occurred. Selection into migration may produce lower or higher vulnerability before the observed shock. For this reason, the manuscript uses the language of association, conditioning and moderation rather than protection.

\section{Data}

Table 1 summarizes the data, samples and variable definitions used in the main analyses. The Kyrgyzstan analysis uses the 2019 wave of the Life in Kyrgyzstan Study. The survey is produced by the LiK consortium and archived through the Research Data Center of IZA \parencite{LIK2023}. The relevant observation unit for the preferred model is the adult respondent, linked to household-level remittance and shock measures. The final model sample is 6,297 adults from 2,215 households. The food-insecurity reference period is 12 months.

The Uzbekistan analysis uses Listening to the Citizens of Uzbekistan rounds 49-82. The producer is the World Bank, and the survey is part of the High-Frequency Phone Surveys collection \parencite{WorldBank2025a}. The household-level data file contains repeated household-round records through round 82. The analysis uses household-rounds as the observation unit and includes 47,135 household-rounds from 2,000 households in the preferred broad-shock model. The food-insecurity reference period is 30 days.

The sample construction rule in both countries is conservative. Records are retained only when the food-insecurity outcome, remittance indicator, shock indicator and required controls are valid. Missing or unresolved codes are not converted into substantive responses. This complete-case approach protects measurement validity but may introduce selection if missingness is related to shocks, remittances or food insecurity.

Kyrgyzstan estimates are unweighted because no approved survey weight is assigned for the preferred analysis. Uzbekistan estimates are unweighted because the available \texttt{popw} documentation remains insufficient for approval in the current design. Standard errors are clustered by household in both countries. Kyrgyzstan uses adult respondent records clustered by household; Uzbekistan uses repeated household-rounds clustered by household.

Kazakhstan is included only as benchmark context. The Kazakhstan data provide adult respondent-year FIES benchmark records for 2014-2017 with year-specific original weights. Because verified remittance and household-shock mechanism variables are absent, Kazakhstan is not used in the remittance-shock interaction analysis.

The country-specific observation units have practical implications. In Kyrgyzstan, multiple adults can appear within the same household, so the food-insecurity outcome has an individual-reporting component while remittances and shocks are household-level exposures. This structure is substantively useful because adult respondents may differ in reported experiences, but it also means standard errors must recognize household clustering. In Uzbekistan, repeated household-round records create a different dependence structure. A household can contribute many observations over time, and its remittance and shock status may change across rounds. The household cluster is therefore the natural unit for inference.

The difference in recall periods is equally important. Kyrgyzstan's 12-month food-insecurity questions may capture hardship that occurred before, during or after a reported household shock. Uzbekistan's 30-day food-insecurity questions are more temporally proximate to the monthly panel structure, although the exact ordering of remittances and shocks within a round may still be unresolved. The paper therefore avoids direct comparison of food-insecurity levels across countries and focuses instead on within-country moderation patterns.

The weighting decision is part of the data design rather than a convenience choice. For Uzbekistan, the L2CU documentation confirms the survey's repeated panel structure and the existence of household data through round 82, but the project did not approve \texttt{popw} for model estimation. Using a weight without sufficient documentation could create a misleading appearance of official population inference. The safer approach is to report unweighted estimates clearly and to identify weighting as a limitation for future revision if documentation becomes available.

The treatment of Kazakhstan follows the same evidentiary rule. Kazakhstan benchmark files contain FIES information and weights, but not the remittance and household-shock mechanism variables required for the research question. Including Kazakhstan in the regression design would require either inventing variables or changing the question. The manuscript does neither. It uses Kazakhstan to show that regional food-insecurity benchmarking is possible while reserving the mechanism analysis for Kyrgyzstan and Uzbekistan.

\section{Measures}

Food insecurity is measured as a raw count from eight FIES-style items. Higher values indicate more food-insecurity experiences. The score is constructed only when all eight items are valid. The raw scores are appropriate for within-country regression and prediction, but they are not official calibrated national prevalence estimates. The different Kyrgyzstan and Uzbekistan reference periods mean that levels should not be read as directly comparable across countries.

Remittance receipt is measured as a household-level indicator. In Kyrgyzstan, remittance receipt is linked from household information to adult respondents. In Uzbekistan, remittance receipt combines verified household remittance components from the L2CU household records. Receipt is preferred over amount because unresolved amount units, timing and missingness create avoidable measurement risk for the primary design.

Household shock exposure is measured country by country. In Kyrgyzstan, the main exposure is any verified household shock, with economic, health, agricultural and climate categories used for robustness. In Uzbekistan, the preferred broad-shock exposure includes household work loss, major illness, major injury and death. Service disruption is not included in the preferred shock variable because it is not a defensible household shock exposure for the main interaction. The work-loss-only Uzbekistan model is retained as secondary event-specific exploratory evidence because its joint remittance-plus-work-loss cell contains only 10 household-rounds from nine households.

Controls include verified demographic and household-composition variables, location variables and time structure available in each country. Current income, asset sales and coping variables are not core controls because they may be downstream of shocks or remittances. Their inclusion could obscure the association the paper is designed to describe. This choice is consistent with the conceptual model rather than a claim that omitted pathways are irrelevant.

\section{Empirical strategy}

The main country-specific model is:

\begin{equation}
FI_{ict}
=
\beta_0
+
\beta_1 Remit_{ict}
+
\beta_2 Shock_{ict}
+
\beta_3
\left(
Remit_{ict}\times Shock_{ict}
\right)
+
X_{ict}'\gamma
+
\delta_c
+
\tau_t
+
\varepsilon_{ict}.
\label{eq:main}
\end{equation}

Here, \(FI_{ict}\) is the food-insecurity raw score for observation \(i\) in country-specific location or household context \(c\) and time \(t\). \(Remit_{ict}\) is remittance receipt, \(Shock_{ict}\) is verified household shock exposure, and \(Remit_{ict} \times Shock_{ict}\) is the interaction of interest. \(X_{ict}\) contains verified controls. \(\delta_c\) denotes location fixed effects where used, and \(\tau_t\) denotes time or round fixed effects where relevant. The error term is clustered by household.

The coefficient \(\beta_3\) is central because it describes whether the shock association differs by remittance status. Since the outcome is higher-worse, a negative \(\beta_3\) is consistent with a weaker adverse shock association among remittance-receiving households. The model is not pooled across countries because LiK and L2CU differ in observation unit, reference period, field design and shock measurement. Pooling would add size but reduce interpretability.

Adjusted predictions are reported for the four remittance-shock groups: no remittance/no shock, remittance/no shock, no remittance/shock and remittance/shock. These predictions translate the interaction into the outcome scale. Interaction contrasts are also reported: the shock association without remittances, the shock association with remittances, the remittance association without shock and the remittance association with shock.

Kyrgyzstan models use household-clustered standard errors, demographic controls, residence and region fixed effects. Uzbekistan preferred models use household-clustered standard errors, household controls and round fixed effects. The Uzbekistan household fixed-effects specification is a qualification that compares within-household variation over rounds. It does not resolve all endogeneity because time-varying confounding, timing of remittances and measurement error remain possible.

Confidence intervals are calculated from clustered standard errors. The primary analysis was conducted in the project Python environment documented in the reproducibility records: Python 3.12.13 with pandas 3.0.1, numpy 2.5.1 and related data-processing packages. The supervisor revision does not estimate new models; it reports the frozen approved results.

The control strategy is intentionally restrained. A common temptation in household welfare models is to include every available socioeconomic variable. In this setting, that approach could be misleading because some variables may be part of the pathway through which shocks and remittances are associated with food insecurity. Current labor income, asset sales, borrowing, reduced spending and coping behavior may all respond to shocks or remittance receipt. Treating them as ordinary pre-treatment controls could absorb meaningful variation and change the interpretation of the interaction. The preferred models therefore use verified demographic, household-composition, location and time variables as the core adjustment set.

The four-group prediction framework is also central to transparency. An interaction coefficient states whether two differences differ, but readers need to see the underlying pattern. The shock contrast among non-remittance households shows the association of shocks in the absence of remittance receipt. The shock contrast among remittance households shows whether that association remains similar when remittances are present. The remittance contrast among non-shocked households separates baseline remittance differences from shock-period differences. The remittance contrast among shocked households shows how remittance status is associated with food insecurity when shocks are observed. Together, these contrasts prevent the discussion from relying on a single product term.

The bounded-outcome robustness check serves a similar interpretive purpose. A raw score bounded between 0 and 8 can be modeled linearly for transparency, but predicted values and uncertainty should be checked against the outcome's limits. Poisson specifications provide one way to respect the count nature of the outcome. The manuscript reports their standardized predictions rather than interpreting the nonlinear interaction coefficient by itself, because nonlinear interactions are not generally readable as simple differences in the same way as linear product terms.

\section{Results}

\subsection{Samples and descriptive patterns}

The preferred Kyrgyzstan model uses 6,297 adult respondents from 2,215 households. Four descriptive groups are available: 4,131 adults with no remittance and no shock, 760 with remittance and no shock, 1,106 with no remittance and shock, and 318 with remittance and shock. All Kyrgyzstan four-group cells are classified as adequate.

The preferred Uzbekistan broad-shock model uses 47,135 household-rounds from 2,000 households. The broad-shock groups include 43,329 household-rounds with no remittance and no verified shock, 3,178 with remittance and no verified shock, 586 with no remittance and verified shock, and 42 with remittance and verified shock. The joint remittance-plus-verified-shock group contains 38 households.

\subsection{Kyrgyzstan preferred estimates}

In Kyrgyzstan, the preferred interaction estimate is -0.2140 with a clustered standard error of 0.2250, a 95\% confidence interval from -0.6549 to 0.2269 and a p-value of 0.3415. The estimate is negative, but the interval includes zero. The approved classification is therefore directional but imprecise. The adjusted predictions are 1.240 for no remittance/no shock, 0.983 for remittance/no shock, 1.449 for no remittance/shock and 0.978 for remittance/shock.

\subsection{Kyrgyzstan robustness}

Kyrgyzstan robustness checks do not replace the preferred model. The standardized interaction is -0.091 with a 95\% confidence interval from -0.278 to 0.096. The bounded-outcome Poisson check gives four-group predictions that follow the same broad pattern: 1.243, 1.033, 1.468 and 1.012. Secondary shock-category models are mixed and imprecise. These checks support a cautious interpretation: the estimated direction is compatible with a weaker shock association among remittance-receiving households, but precision is limited.

\subsection{Uzbekistan broad-shock estimates}

The preferred Uzbekistan broad-shock model estimates an interaction of -0.5406 with a clustered standard error of 0.2555, a 95\% confidence interval from -1.0415 to -0.0398 and a p-value of 0.03437. The shock association without remittances is 0.5538, while the shock association with remittances is 0.0132. The remittance association without shock is -0.1382, and the remittance association with shock is -0.6788. These contrasts show that the broad-shock association with food insecurity is concentrated among non-remittance household-rounds in the preferred model.

\subsection{Uzbekistan adjusted predictions}

Uzbekistan adjusted predictions show the same pattern on the food-insecurity scale. Predicted scores are 0.734 for no remittance/no verified shock, 0.596 for remittance/no verified shock, 1.288 for no remittance/verified shock and 0.609 for remittance/verified shock. The remittance-plus-shock prediction has wide uncertainty because the joint group is modest, but the four-group pattern is central to the interpretation. The preferred Uzbekistan evidence is classified as a moderate conditional association with limitations.

\subsection{Uzbekistan fixed-effects qualification}

The household fixed-effects model is a prominent qualification. The fixed-effects interaction is -0.1771 with a clustered standard error of 0.1910, a 95\% confidence interval from -0.5515 to 0.1973 and a p-value of 0.3539. The estimate remains negative but is attenuated and imprecise. The fixed-effects sample includes 47,135 observations from 2,000 households, with 476 remittance switchers, 429 shock switchers and 130 households switching both. This qualification suggests that stable household differences and limited within-household switching are important for interpretation.

\subsection{Temporal and participation sensitivities}

Temporal and participation sensitivities from Phase 7 are used as qualifications rather than new primary findings. They support the need to treat the Uzbekistan result as a conditional association rather than a fully robust causal parameter. They also reinforce the importance of round fixed effects, household clustering and transparent reporting of L2CU participation patterns.

\subsection{Work-loss exploratory result}

The Uzbekistan work-loss-only result is secondary event-specific exploratory evidence. Its interaction estimate is negative, but the joint remittance-plus-work-loss cell contains only 10 household-rounds from nine households. That cell is too sparse to support the main Uzbekistan conclusion. The broad verified-shock result is therefore the preferred Uzbekistan model, while work loss is retained as an appendix-level warning about event-specific rarity.

\subsection{Standardized directional comparison}

The standardized interaction is -0.091 for Kyrgyzstan and -0.337 for Uzbekistan. This comparison is directional only. It does not pool countries and does not imply that the surveys share the same outcome timing, shock definition or observation unit. The comparison supports the statement that both preferred estimates are negative, while the Uzbekistan estimate is larger in standardized magnitude and more precise.

\subsection{Kazakhstan benchmark}

Kazakhstan is reduced to benchmark context in the main text. The validated annual benchmark file reports weighted mean raw scores of 0.802 in 2014, 0.528 in 2015, 0.680 in 2016 and 0.821 in 2017, using year-specific original weights. These are benchmark summaries from supplied FIES variables, not official calibrated national prevalence estimates. Because Kazakhstan lacks verified remittance and shock mechanism variables, it is not part of the regression design.

The descriptive cells help readers understand the precision of the estimates. Kyrgyzstan has enough observations in each of the four main groups, but the remittance-plus-shock group is still much smaller than the no-remittance/no-shock group. Uzbekistan's broad-shock joint group is classified as usable, yet it contains only 42 household-rounds from 38 households. This imbalance is normal in interaction designs with relatively rare combined exposure, but it affects how much weight should be placed on a single estimate. It is the reason the paper pairs the preferred model with fixed-effects, influence and sensitivity checks.

The Uzbekistan adjusted predictions are particularly informative. The predicted score for non-remittance household-rounds rises sharply when a verified shock is observed. Among remittance household-rounds, the corresponding difference is close to zero in the preferred model. This is the substantive pattern behind the negative interaction. The wide interval around the remittance-plus-shock prediction signals uncertainty, but the contrast framework makes the shape of the association clear.

For Kyrgyzstan, the adjusted predictions point in the same conceptual direction, but the uncertainty is greater. The remittance-plus-shock group does not show a higher predicted score than the remittance/no-shock group in the preferred model, while the no-remittance shock group has a higher predicted score than the no-remittance/no-shock group. However, the interaction interval includes zero. The correct interpretation is therefore not absence of a pattern, but insufficient precision for a strong conclusion.

The Kazakhstan benchmark is intentionally compact in the results because it does not answer the research question. The annual benchmark values are useful for regional context and for documenting that Kazakhstan files were not ignored. They should not distract from the two-country empirical design. The appendix and benchmark table retain the validated values for readers who want them.

\section{Discussion}

The main finding is not that remittances provided demonstrated protection. The more credible interpretation is that the observed relationship between shocks and food insecurity differs by remittance status, with stronger evidence in Uzbekistan than in Kyrgyzstan. This interpretation fits the informal-insurance literature but remains bounded by the observational design.

The Kyrgyzstan estimate points in the expected direction but is imprecise. There are several possible reasons. The Kyrgyzstan outcome is measured at the adult respondent level and uses a 12-month reference period. Household-level shocks and remittances are linked to individual reports, and variation within households may add noise. The final joint remittance-plus-shock cell is adequate, but the confidence interval remains wide. Secondary shock-category models also do not produce a uniform precise pattern. Kyrgyzstan therefore contributes directional evidence, not a strong standalone claim.

The Uzbekistan broad-shock model is clearer in the preferred pooled specification. The four-group predictions show that verified shocks are associated with a large increase in predicted food-insecurity scores among non-remittance household-rounds, while the shock association among remittance household-rounds is near zero. This pattern is consistent with the idea that remittance-receiving households may have resources or networks that coincide with lower food-insecurity vulnerability when shocks occur.

The fixed-effects qualification prevents overstatement. Once household fixed effects are included, the Uzbekistan interaction remains negative but becomes smaller and imprecise. This may mean that part of the preferred-model pattern reflects stable household differences rather than within-household changes. It may also reflect limited switching in remittance and shock status or insufficient information in the joint switching group. The fixed-effects result should therefore appear in any summary of the Uzbekistan findings.

The results align with \textcite{ChoiYang2007}'s remittance-insurance evidence and with the broader consumption-stabilization literature \parencite{CombesEbeke2011}, but the present design is more modest. Choi and Yang use rainfall shocks as instruments for income changes in the Philippines. This paper does not have an equivalent identification strategy. It instead reports a descriptive interaction using verified country-specific household survey measures. That difference is central to the level of claims.

The Kyrgyzstan result should also be interpreted alongside LiK-based remittance-spending work. \textcite{Wang2021} show that remittances are related to household budget allocation and that endogeneity and heterogeneity matter. The present Kyrgyzstan model similarly finds a negative moderation pattern but with limited precision. Rather than contradicting prior LiK research, this result suggests that food-insecurity moderation is a harder and noisier outcome to estimate with the available 2019 food-insecurity module.

The Uzbekistan result should be interpreted alongside \textcite{Kakhkharov2021}, who show that remittance-receiving households in Uzbekistan differ in expenditure allocation and explicitly address selection. The present analysis does not reproduce their expenditure model; it adds a food-insecurity and shock-moderation question using L2CU. The shared lesson is that Uzbekistan remittance evidence must be read with selection in mind.

Cross-country comparison is limited by design. Kyrgyzstan uses adult respondents and a 12-month food-insecurity reference period; Uzbekistan uses household-rounds and a 30-day reference period. Shock measures differ, and the surveys cover different periods. The paper therefore compares the direction and credibility of country-specific moderation patterns rather than food-insecurity levels or pooled coefficients.

The plausible mechanism is liquidity and informal support, but it is not directly demonstrated. Remittances may allow households to maintain food purchases after shocks, or they may be sent in response to a shock. They may also proxy for household networks, prior wealth, migrant employment opportunities or unobserved resilience. The adjusted models and fixed-effects check narrow some concerns but do not remove all alternative explanations.

Overall, the evidence is useful because it identifies heterogeneity in the shock-food-insecurity association. Households without remittances appear more exposed to the adverse association in the Uzbekistan preferred model, while Kyrgyzstan shows a similar direction with greater uncertainty. That pattern is relevant for resilience analysis, but it should guide questions for further evaluation rather than direct claims about what remittance policy would do.

One way to read the findings is through the distinction between exposure and capacity. Shocks mark exposure to adverse events; remittance receipt marks one possible dimension of household capacity or network access. The study finds that the association between exposure and food-insecurity scores is weaker among remittance recipients in the Uzbekistan broad-shock model and points in the same direction in Kyrgyzstan. That is an important descriptive pattern, but capacity is not fully observed. Households may have savings, livestock, social support, public transfers or informal credit that are not captured by the remittance indicator.

The fixed-effects result sharpens the interpretation because it shifts attention from between-household differences to within-household changes. If remittance-receiving households are systematically different even before shocks occur, the pooled model can partly reflect those differences. The fixed-effects model reduces reliance on stable between-household differences, and the Uzbekistan estimate becomes smaller and imprecise. This does not invalidate the pooled association, but it changes what can be said. The safest statement is that Uzbekistan shows a moderate conditional association in the preferred model, qualified by attenuated and imprecise within-household evidence.

The result also has implications for future data collection. Stronger evidence would require clearer timing of shocks, remittance receipt, transfer amounts and food-insecurity experiences. A panel module that records when a shock occurred, when remittances arrived, whether transfers changed because of the shock and how households adjusted food consumption would allow a more direct mechanism test. The current data support a moderation analysis, not a full pathway decomposition.

\section{Policy implications}

The findings have policy relevance for vulnerability assessment, not for claiming that any policy instrument has been evaluated. In the Uzbekistan preferred model, verified shocks are associated with substantially higher food-insecurity scores among non-remittance household-rounds, while the same association is much weaker among remittance household-rounds. This suggests that information on household support networks may help identify which shocked households are most vulnerable. It does not show that increasing remittances would lower food insecurity.

Shock-responsive social protection frameworks emphasize the need for systems that can identify shocks, reach affected households and adapt support when risks materialize \parencite{OBrien2018,Bowen2020}. The present results fit that agenda by highlighting heterogeneity in observed hardship after shocks. Households without remittance support may deserve particular attention in vulnerability screening, especially when employment, health or death shocks occur.

Private remittances cannot replace public social protection. Transfers from migrants are uneven, uncertain and shaped by migrant labor-market conditions. Remittance-receiving households may still experience food insecurity, and households without migrants may lack access to this informal support. The policy implication is therefore not to rely on remittances, but to recognize them as one part of the household resource environment.

The findings also suggest a research agenda for program design. If governments or development partners want to evaluate lower transfer costs, cash assistance, social insurance or employment support, those interventions require separate evaluation designs. This study can motivate such evaluations by showing where vulnerability appears concentrated, but it does not estimate the result of any program.

Kazakhstan's role is only contextual. Benchmark FIES summaries can help situate Central Asian food-security monitoring, but they do not inform remittance-shock policy mechanisms in the current project. Future Kazakhstan microdata with verified remittance and shock variables would be needed before extending the interaction design.

The policy value of the paper is therefore diagnostic. It identifies a subgroup distinction that may matter for targeting and monitoring: shocked households with and without remittance support. In a shock-responsive system, such information could be combined with poverty, household composition, location and exposure indicators. It should not be used mechanically. Some remittance households may still be vulnerable, and some non-remittance households may have other support. The point is to improve the questions asked by vulnerability systems, not to create a single eligibility rule.

The paper also suggests that migration and social protection should be analyzed together rather than in separate silos. Migration policy affects households, but social-protection systems remain responsible for public risk management. A household's private transfer network may reduce some observed hardship, but it is uneven and uncertain. This is consistent with adaptive social-protection frameworks that treat resilience as a systems issue involving information, delivery capacity and timely support.

\section{Limitations}

The study is observational. Remittance receipt is selected, and migrant-sending households may differ from other households before observed shocks occur. Reverse timing is also possible: remittances may be sent in response to shocks rather than being available beforehand. Residual confounding remains possible even after controls, fixed effects and household clustering.

The Uzbekistan household fixed-effects result is a major limitation and qualification. It remains negative but is attenuated and imprecise, which means the preferred pooled association should not be read as a stable within-household estimate. The fixed-effects result may reflect stable household selection, limited switching or reduced information in the joint switching group.

Weighting is another limitation. L2CU estimates are unweighted because \texttt{popw} documentation was not approved for this analysis, and the Kyrgyzstan models do not use an approved LiK survey weight. The results therefore describe the analyzed samples and model structure rather than approved weighted population estimates.

Measurement differences also matter. Kyrgyzstan uses adult respondent outcomes with a 12-month reference period, while Uzbekistan uses household-round outcomes with a 30-day reference period. Shock definitions differ, and the Uzbekistan broad-shock joint cell, though usable, is modest. The work-loss-only joint cell is sparse and therefore exploratory. Complete-case analysis may introduce selection if missingness is related to food insecurity, remittances or shocks. Measurement error is possible in self-reported food-insecurity items, remittance indicators and shock reports. Kazakhstan remains benchmark-only because mechanism variables are absent.

A final limitation concerns interpretation of null or imprecise results. Imprecision does not mean the relationship is absent, and a negative estimate does not mean the mechanism is proven. The Kyrgyzstan result sits exactly in that space: the sign is substantively meaningful, but the confidence interval is wide. The Uzbekistan fixed-effects estimate also sits in that space: it is directionally consistent, but too imprecise to carry the preferred conclusion on its own. The manuscript therefore uses evidence classifications rather than a binary significant/not-significant narrative.

\section{Conclusion}

This paper examined whether the positive association between household shocks and food insecurity is weaker among remittance-receiving households in Kyrgyzstan and Uzbekistan. The results are country-specific and observational. Kyrgyzstan's preferred interaction is negative but imprecise. Uzbekistan's preferred broad-shock interaction is negative and more precise, but the household fixed-effects estimate is attenuated and imprecise. The work-loss-only Uzbekistan evidence remains secondary and exploratory because of a sparse joint cell.

The best summary is therefore balanced. The evidence is directionally consistent with remittance-related moderation in both countries, with stronger support in the Uzbekistan broad-shock model and a prominent fixed-effects qualification. The findings do not establish that remittances provided household protection, and they do not justify replacing public social protection with private transfers. They do support continued attention to remittances, shocks and food insecurity as linked dimensions of household resilience in Central Asia.

\section{Declarations}

\subsection*{Funding}
This research received no external funding.

\subsection*{Conflicts of interest}
The author declares no conflict of interest.

\subsection*{Ethics statement}
This study used de-identified secondary survey data. No new participants were recruited, and the author did not access directly identifying personal information. The data were analyzed in accordance with the access and use conditions established by the original data providers. The original survey organizations were responsible for participant consent and applicable ethical procedures.

\subsection*{Data availability}
The source datasets are available from their respective data providers subject to registration, licensing, and access conditions. The author does not own the underlying microdata and therefore cannot redistribute them with this article. Non-disclosive derived outputs and replication materials will be made available in a public repository before publication.

\subsection*{Code availability}
Analysis code and non-restricted replication materials are available from the corresponding author upon reasonable request and will be deposited in a public repository before publication.

\subsection*{Author contributions}
Foziljon Alisherov: Conceptualization, methodology, software, validation, formal analysis, investigation, data curation, visualization, writing---original draft, writing---review and editing, and project administration.

\subsection*{Acknowledgements}

The author gratefully acknowledges Sarvinoz Abdumo'minova, Ibroxim Ergashev,
Nilufar Farmonova, Ozodbek Islomov, and Jahongir Boltayev for their assistance
with preliminary data organization and literature searches.

\subsection*{Generative-AI disclosure}
Generative AI tools were used to assist with drafting, organization, and language refinement. The author reviewed and verified all analyses, numerical results, citations, interpretations, and final text and takes full responsibility for the content.

\printbibliography[title={References}]

\clearpage

\appendix
\section{Survey and variable construction}
This appendix records approved supplementary material for the author-review package. Kyrgyzstan uses the 2019 Life in Kyrgyzstan Study adult respondent analysis file. Uzbekistan uses Listening to the Citizens of Uzbekistan rounds 49--82 at the household-round level. Kazakhstan is benchmark only.

\section{Sample flow and missingness}
The preferred Kyrgyzstan model uses 6,297 adult respondents from 2,215 households. The preferred Uzbekistan broad-shock model uses 47,135 household-rounds from 2,000 households. Complete-case restrictions are documented in the Phase 8 manuscript and prior checkpoints.

\section{Food-insecurity measurement checks}
Food insecurity is measured with eight FIES-style items summed to a raw score from 0 to 8. Raw scores are not official calibrated national prevalence estimates.

\section{Full regression results}
The approved primary interactions are KG\_M2: -0.2140 (95\% CI -0.6549 to 0.2269; p=0.3415) and UZBROAD\_M2: -0.5406 (95\% CI -1.0415 to -0.0398; p=0.03437). No new models are estimated in Phase 9.

\section{Bounded-outcome robustness}
Bounded-outcome robustness is interpreted through standardized four-group predictions rather than the raw nonlinear interaction coefficient alone.

\section{Uzbekistan household fixed effects}
The household fixed-effects interaction is -0.1771 (95\% CI -0.5515 to 0.1973; p=0.3539). It is directionally consistent but attenuated and imprecise.

\section{Lagged and participation sensitivity}
Temporal and participation sensitivities are retained as qualifications, not replacements for frozen primary findings.

\section{Influence diagnostics}
Rare-cell influence checks are documented in the Phase 5 and Phase 7 checkpoints. Uzbekistan's broad remittance-plus-shock group contains 42 household-rounds from 38 households.

\section{Work-loss exploratory analysis}
\textbf{SECONDARY EVENT-SPECIFIC EXPLORATORY ANALYSIS.} The remittance-plus-work-loss group contains 10 household-round observations from nine households. This result does not replace the approved Uzbekistan broad-shock model.

\section{Kazakhstan benchmark methodology}
Kazakhstan 2014--2017 FIES files are used for regional benchmark context only. The benchmark uses year-specific original weights and supplied FIES probability variables; it is not a remittance-shock mechanism analysis.


\clearpage

\begin{table}[!htbp]
\centering
\caption{Data, samples and variable definitions}
\label{tab:data}
\footnotesize
\setlength{\tabcolsep}{4pt}
\begin{tabularx}{\linewidth}{p{0.14\linewidth} X X X}
\toprule
Country & Survey and sample & Main variables & Design notes \\
\midrule
Kyrgyzstan & Life in Kyrgyzstan Study, 2019 wave. Final sample: 6,297 adult respondents from 2,215 households. & Outcome: eight-item food-insecurity raw score, 12-month reference period. Remittance: household receipt. Shock: any verified household shock. & Household-clustered standard errors; demographic controls; residence and region fixed effects; unweighted because no approved analysis weight is assigned. \\
Uzbekistan & Listening to the Citizens of Uzbekistan Survey, rounds 49--82. Final sample: 47,135 household-rounds from 2,000 households. & Outcome: eight-item food-insecurity raw score, 30-day reference period. Remittance: verified household receipt. Shock: work loss, major illness, major injury or death. & Household-clustered standard errors; household controls; round fixed effects; unweighted because \texttt{popw} is not approved. \\
Kazakhstan & FIES benchmark files, 2014--2017 adult respondent-years. & FIES benchmark outcomes only. Remittance and household-shock mechanism variables are not available. & Benchmark only; year-specific original weights; not part of the regression design. \\
\bottomrule
\end{tabularx}
\begin{minipage}{0.95\linewidth}\footnotesize Notes: Countries are not pooled. Internal source-variable names are reserved for appendices and project code.\end{minipage}
\end{table}

\clearpage

\begin{table}[!htbp]
\centering
\caption{Four-group adjusted predictions}
\label{tab:fourgroups}
\footnotesize
\setlength{\tabcolsep}{5pt}
\begin{tabularx}{\linewidth}{p{0.15\linewidth} p{0.16\linewidth} X p{0.16\linewidth} p{0.18\linewidth}}
\toprule
Country & Model & Group & Prediction & Observations / households \\
\midrule
Kyrgyzstan & KG\_M2 & No remittance, no shock & 1.240 & 4,131 / 1,447 \\
Kyrgyzstan & KG\_M2 & Remittance, no shock & 0.983 & 760 / 281 \\
Kyrgyzstan & KG\_M2 & No remittance, shock & 1.449 & 1,106 / 381 \\
Kyrgyzstan & KG\_M2 & Remittance, shock & 0.978 & 318 / 112 \\
Uzbekistan & UZBROAD\_M2 & No remittance, no verified shock & 0.734 [0.677, 0.792] & 43,329 / 1,990 \\
Uzbekistan & UZBROAD\_M2 & Remittance, no verified shock & 0.596 [0.499, 0.694] & 3,178 / 477 \\
Uzbekistan & UZBROAD\_M2 & No remittance, verified shock & 1.288 [1.094, 1.483] & 586 / 407 \\
Uzbekistan & UZBROAD\_M2 & Remittance, verified shock & 0.609 [0.133, 1.085] & 42 / 38 \\
\bottomrule
\end{tabularx}
\begin{minipage}{0.95\linewidth}\footnotesize Notes: Brackets show 95\% confidence intervals where available in the frozen register. Estimates are unweighted; Uzbekistan \texttt{popw} is not approved.\end{minipage}
\end{table}

\clearpage

\begin{table}[!htbp]
\centering
\caption{Kyrgyzstan preferred model}
\label{tab:kgmodel}
\footnotesize
\begin{tabularx}{\linewidth}{X p{0.14\linewidth} p{0.14\linewidth} p{0.23\linewidth} p{0.12\linewidth}}
\toprule
Term & Estimate & Clustered SE & 95\% CI & p-value \\
\midrule
Remittance $\times$ shock interaction & -0.2140 & 0.2250 & -0.6549 to 0.2269 & 0.3415 \\
\bottomrule
\end{tabularx}
\begin{minipage}{0.95\linewidth}\footnotesize Notes: Model KG\_M2; 6,297 adult respondents; 2,215 household clusters; preferred adjusted controls; region/residence fixed effects; unweighted. Evidence classification: directional but imprecise. No significance stars are used.\end{minipage}
\end{table}

\clearpage

\begin{table}[!htbp]
\centering
\caption{Uzbekistan broad-shock preferred model and fixed-effects qualification}
\label{tab:uzmodel}
\footnotesize
\begin{tabularx}{\linewidth}{X p{0.14\linewidth} p{0.14\linewidth} p{0.23\linewidth} p{0.12\linewidth}}
\toprule
Specification & Estimate & Clustered SE & 95\% CI & p-value \\
\midrule
Preferred broad-shock interaction, UZBROAD\_M2 & -0.5406 & 0.2555 & -1.0415 to -0.0398 & 0.03437 \\
Household fixed-effects interaction, UZBROAD\_FE\_HH & -0.1771 & 0.1910 & -0.5515 to 0.1973 & 0.3539 \\
\bottomrule
\end{tabularx}
\begin{minipage}{0.95\linewidth}\footnotesize Notes: 47,135 household-rounds; 2,000 household clusters; household-clustered standard errors; preferred model uses verified household-composition controls and round fixed effects. Fixed-effects specification includes household and round fixed effects. Estimates are unweighted because \texttt{popw} is not approved. No significance stars are used.\end{minipage}
\end{table}

\clearpage

\begin{table}[!htbp]
\centering
\caption{Interaction contrasts and robustness qualifications}
\label{tab:contrasts}
\footnotesize
\setlength{\tabcolsep}{5pt}
\begin{tabularx}{\linewidth}{p{0.18\linewidth} X p{0.13\linewidth} p{0.22\linewidth} p{0.12\linewidth}}
\toprule
Country/model & Contrast or qualification & Estimate & 95\% CI & p-value \\
\midrule
Kyrgyzstan KG\_M2 & Shock association without remittances & 0.2094 & -0.0315 to 0.4502 & 0.0884 \\
Kyrgyzstan KG\_M2 & Shock association with remittances & -0.0047 & -0.3800 to 0.3707 & 0.9806 \\
Uzbekistan UZBROAD\_M2 & Shock association without remittances & 0.5538 & 0.3689 to 0.7388 & $<$0.001 \\
Uzbekistan UZBROAD\_M2 & Shock association with remittances & 0.0132 & -0.4541 to 0.4805 & 0.9559 \\
Uzbekistan UZBROAD\_FE\_HH & Fixed-effects qualification & -0.1771 & -0.5515 to 0.1973 & 0.3539 \\
Uzbekistan work-loss only & Sparse-cell exploratory warning & -1.1119 & -1.8088 to -0.4151 & 0.0018 \\
\bottomrule
\end{tabularx}
\begin{minipage}{0.95\linewidth}\footnotesize Notes: Work-loss-only result is secondary event-specific exploratory analysis; the remittance-plus-work-loss group contains 10 household-round observations from nine households. No significance stars are used.\end{minipage}
\end{table}

\clearpage

\begin{table}[!htbp]
\centering
\caption{Standardized directional comparison}
\label{tab:standardized}
\footnotesize
\begin{tabularx}{\linewidth}{p{0.15\linewidth} p{0.18\linewidth} p{0.18\linewidth} p{0.18\linewidth} X}
\toprule
Country & Model & Standardized interaction & 95\% CI & Interpretation \\
\midrule
Kyrgyzstan & KG\_R\_STD & -0.091 & -0.278 to 0.096 & Directional but imprecise \\
Uzbekistan & UZBROAD standardized & -0.337 & -0.649 to -0.025 & Moderate conditional association with fixed-effects qualification \\
\bottomrule
\end{tabularx}
\begin{minipage}{0.95\linewidth}\footnotesize Notes: Not a pooled comparison. Kyrgyzstan has 6,297 observations and 2,215 clusters; Uzbekistan has 47,135 observations and 2,000 clusters. Both are unweighted; Uzbekistan \texttt{popw} is not approved.\end{minipage}
\end{table}

\clearpage

\begin{table}[!htbp]
\centering
\caption{Kazakhstan benchmark}
\label{tab:kazakhstan}
\footnotesize
\begin{tabularx}{\linewidth}{p{0.14\linewidth} p{0.18\linewidth} p{0.18\linewidth} p{0.20\linewidth} X}
\toprule
Year & Eligible observations & Weighted mean raw score & Mean supplied moderate-or-severe probability & Interpretation \\
\midrule
2014 & 898 & 0.802 & 0.082 & Benchmark only \\
2015 & 926 & 0.528 & 0.048 & Benchmark only \\
2016 & 936 & 0.680 & 0.078 & Benchmark only \\
2017 & 968 & 0.821 & 0.093 & Benchmark only \\
\bottomrule
\end{tabularx}
\begin{minipage}{0.95\linewidth}\footnotesize Notes: Uses year-specific original weights and supplied FIES variables. Values are benchmark summaries, not a remittance-shock mechanism test.\end{minipage}
\end{table}

\clearpage

\begin{figure}[!htbp]
\centering
\fbox{\begin{minipage}{0.86\linewidth}
\centering
\textbf{Household shocks}\\[0.35em]
$\Downarrow$ \quad food-insecurity risk \quad $\Downarrow$\\[0.35em]
\textbf{Food-insecurity raw score}\\[0.7em]
\hrule
\vspace{0.6em}
\textbf{Remittance receipt as observed moderator}\\
The empirical design tests whether the shock--food-insecurity association differs by remittance status. The framework is non-causal.
\end{minipage}}
\caption{Conceptual framework: remittances as an observational moderator of household shocks and food insecurity.}
\label{fig:conceptual}
\end{figure}

\clearpage

\begin{figure}[!htbp]
\centering
\includegraphics[width=0.92\linewidth,height=0.72\textheight,keepaspectratio]{\detokenize{../../outputs/figures/figure_19_kyrgyzstan_adjusted_four_groups_v2.png}}
\caption{Kyrgyzstan adjusted food-insecurity predictions by remittance and shock status. Model KG\_M2; adult respondent unit; 6,297 observations; 2,215 household clusters; household-clustered 95\% confidence intervals; unweighted.}
\label{fig:kgpred}
\end{figure}

\clearpage

\begin{figure}[!htbp]
\centering
\includegraphics[width=0.92\linewidth,height=0.72\textheight,keepaspectratio]{\detokenize{../../outputs/figures/figure_25_uzbekistan_broad_shock_predictions.png}}
\caption{Uzbekistan adjusted food-insecurity predictions by remittance and verified-shock status. Model UZBROAD\_M2; household-round unit; 47,135 observations; 2,000 household clusters; household-clustered 95\% confidence intervals; unweighted because \texttt{popw} is not approved.}
\label{fig:uzpred}
\end{figure}

\clearpage

\begin{figure}[!htbp]
\centering
\includegraphics[width=0.92\linewidth,height=0.72\textheight,keepaspectratio]{\detokenize{../../outputs/figures/figure_26_revised_standardized_interactions.png}}
\caption{Standardized remittance--shock interaction associations. Country-specific standardized coefficients with a zero reference line; countries differ in survey unit, recall period and shock definition.}
\label{fig:std}
\end{figure}

\clearpage

\begin{figure}[!htbp]
\centering
\includegraphics[width=0.92\linewidth,height=0.72\textheight,keepaspectratio]{\detokenize{../../outputs/figures/figure_24_kazakhstan_benchmark_with_ci.png}}
\caption{Kazakhstan annual food-insecurity benchmark. Weighted means of supplied probability variables and raw scores; benchmark only, not a mechanism test.}
\label{fig:kaz}
\end{figure}

\end{document}
